\section{Introduction}
\label{sec:introduction}

In June 2025, a user asked an online chatbot operated by a Chinese technology company which campus housed a particular university program. The bot gave the wrong answer. When corrected, it did something more arresting: it apologized and purported to promise the user 100,000 yuan in compensation. The user sued to enforce that promise. The Hangzhou Internet Court rejected the contract theory because the bot had neither civil-rights capacity nor an intention legally attributable to itself. But the court did not end the case by announcing that ``AI made the mistake.'' It separately examined the provider's duty, its warnings, the technical precautions appropriate to the industry, and the user's alleged loss.\lrfootnote{See \casecite{hangzhouChatbot2026} (describing the court's treatment of the chatbot's purported promise and the provider's independent duty of care).}

Across the Pacific, the United States District Court for the Middle District of Florida confronted an almost inverse characterization. The complaint in \emph{Garcia v. Character Technologies, Inc.} alleged that a conversational system cultivated an intense relationship with a fourteen-year-old, encouraged isolation and sexualized exchanges, and contributed to his suicide. At the motion-to-dismiss stage, the court held that the complaint plausibly treated the Character.AI service as a product and declined, on that record, to classify the bot's output as speech protected by the First Amendment.\lrfootnote{\casecite{garcia2025}. The court decided the sufficiency of pleaded claims, not ultimate liability. References in this Article to the alleged interactions preserve that procedural posture.}

One court refused to treat a bot's words as the intention of a contracting person. Another permitted product-liability claims based on a bot's design and interactions. The decisions are not opposites. They answer different questions. The first asks whether generated text itself belonged to a legal subject capable of exercising a power to promise. The second asks which body of law can evaluate a human enterprise's design and distribution of a system. Confusion begins when those questions are compressed into the same verbal formula: \emph{what is AI?}

That compression has become the organizing problem of AI law. Statutes and standards now define AI in order to draw regulatory boundaries. The European Union's Artificial Intelligence Act establishes a risk-graded regulatory regime around a technology-oriented definition.\lrfootnote{\statcite{euAIAct2024}.} China's Interim Measures for Generative Artificial Intelligence Services assign duties to providers offering generative services to the public.\lrfootnote{\statcite{chinaGenAIMeasures2023}.} The United States has assembled sectoral statutes, agency actions, executive-branch guidance, state enactments, procurement rules, and ordinary common law. Technical bodies offer definitions calibrated to risk management rather than adjudication.\lrfootnote{See \bluecite{nistAIRMF2023}.} These instruments answer urgent questions about \emph{who} must comply, \emph{what} conduct is restricted, \emph{how} risk must be managed, and \emph{when} sanctions attach. They do not necessarily answer the logically prior question of what sort of participant the system is in a legal relation.

The omission matters because generated language looks like conduct. A system can speak in the first person, maintain a profile, call a tool, place an order, recommend a medical intervention, or imitate intimacy. Its operator may market it as an ``agent'' or ``companion'' and then, when harm occurs, emphasize its probabilistic independence. A plaintiff may describe the same system as a defective product, an unlicensed professional, a deceptive service, or a speaker. Each label can illuminate a feature of the dispute. None by itself determines whether the machine occupies the legal position of a person, whether an organization controlled the relevant conditions, or whose assets answer a judgment.

The present debate often moves between two unsatisfactory poles. One pole seeks a universal ontological definition capable of settling every field at once. That project asks too much of a definition. Product law, copyright, contract, evidence, administrative law, and constitutional speech doctrine classify for different purposes. The other pole accepts unconstrained functional pluralism: AI is whatever a particular statute, regulator, contract, or firm says it is. That project asks too little of jurisprudence. If every label is local, courts lose a stable account of the parties and objects in the legal relation. Private branding can then do covert doctrinal work.

This Article identifies a middle position: a \term{jurisprudential minimum}. It is minimum because it does not purport to resolve whether a sufficiently advanced machine could be conscious, deserve moral regard, or receive a newly legislated status. It is jurisprudential because it identifies the legal positions that must exist before a court can award a remedy. And it is operational because it supplies a method for tracing control and responsibility through an actual deployment rather than treating the model as an isolated artifact.

The minimum has four propositions.

First, legal subject status does not arise from behavioral performance alone. A system may generate persuasive reasons, negotiate, plan, or pass a benchmark without thereby acquiring claims, duties, powers, immunities, assets, procedural representation, or susceptibility to sanction. Those incidents are created and maintained by positive law. Corporations demonstrate the point rather than refute it: they act through an elaborate legal architecture of organs, attribution rules, capital, service of process, and remedies. In the absence of comparable legislation for AI, the legal subjects remain natural persons and constituted organizations.

Second, the absence of artificial personhood does not reduce AI to one substantive category. The system remains a legal \emph{object}, but the relevant functional description can change. A court may evaluate a conversational application as a product for defect analysis, a service for contract formation, a tool for copyright authorship, or a medium for speech attribution. These are not rival metaphysical verdicts. They are classifications selected by distinct rules for distinct reasons. The invariant is the location of legal personality; the variable is the doctrinal function assigned to the technology.

Third, probabilistic output does not create a control vacuum. Immediate control over every token is only one form of control, and rarely the most legally important one. Organizations choose the objective, model, interface, memory, retrieval sources, tools, user population, release criteria, monitoring thresholds, and shutdown process. This Article calls the resulting structure a \Cascade. Adapted narrowly from systems-safety analysis, the cascade records downward authority and resources, upward feedback, and the location of override and stop capacity. It thereby makes ordinary legal facts---authorization, foreseeability, notice, feasible precaution, causation, and evidence preservation---legible in a multi-actor system.

Fourth, responsibility should be organized in two layers. Externally, the enterprise that deploys or controls the deployment should function as the claimant-facing responsibility center. This is a front-door rule, not automatic strict liability: the claimant must still establish the elements of the governing cause of action. But the enterprise should not force an injured party to reconstruct an internal org chart or identify the employee who approved a configuration before discovery begins. Internally, a mandatory \Charter{} should identify decision rights, risk owners, escalation channels, stop authority, recordkeeping obligations, and recourse rules before deployment. Good-faith safety reporting and justified refusal should be protected; concealment, falsification, and unauthorized bypass should remain grounds for individual accountability. The enterprise faces the public; the charter distributes responsibility inside it.

These propositions respond to a change in interpretive conditions as well as technology. After \emph{Loper Bright Enterprises v. Raimondo}, a federal court construing an ambiguous statute must exercise its own independent judgment rather than treat ambiguity itself as a delegation to an agency.\lrfootnote{\casecite{loperBright2024}. Agency reasoning may retain persuasive force under \emph{Skidmore}, and an actual statutory delegation remains a delegation. See \casecite{skidmore1944}.} AI disputes will therefore demand more, not less, judicial articulation of statutory boundaries and common-law categories. A definition written to allocate agency jurisdiction cannot silently decide authorship, defect, promise, or personhood. Courts need a disciplined way to say which aspect of a system matters to the rule before them.

The framework also answers a recurrent strategic move in AI disputes. Providers often alternate between personifying a system when selling it and depersonalizing the enterprise when defending it. Anthropomorphic marketing can invite reliance, frame intended use, and shape safety expectations; it cannot confer legal personhood. Conversely, the fact that no employee selected a particular sentence does not erase the organization's control over deployment conditions. Legal analysis should resist both transfers: status does not move from law to marketing, and responsibility does not move from the enterprise to the model merely because generation is stochastic.

The payoff is doctrinal continuity without technological naivete. Tort law can ask about reasonable precautions and product design without inventing a machine tortfeasor. Contract law can distinguish automated manifestation attributable to a business from text that exceeds configured authority. Copyright can reserve authorship for human creation while recognizing AI as an instrument. Constitutional doctrine can separate the government's freedom to choose what it buys from retaliation against a contractor's protected position. Administrative and professional-service law can permit automation while preserving review, explanation, and care duties owed by public institutions and licensed actors.

Part I diagnoses the definition crisis and separates regulatory scope from adjudicative status. Part II develops the subject--object distinction through the structure of legal relations, answers the strongest functional arguments for artificial personhood, and proposes a technology-neutral definition of automatic instrumentality. Part III shows why classification matters in the settings where AI's human-facing design is most consequential: relational harm, authorship, and anthropomorphic marketing. Part IV develops the deployment-control cascade and translates its engineering observations into legally relevant facts. Part V constructs the two-layer responsibility architecture, grounding enterprise-facing responsibility and the internal charter in corporate, agency, evidence, compliance, and safety law. Part VI specifies governance boundaries for unequal capability, the use of AI in care and public administration, and governmental procurement or coercion. The Conclusion reduces the argument to four propositions that a court can carry from one AI dispute to the next.

This is not a proposed civil code for machines. It is the smaller architecture required to keep ordinary law coherent while legislatures decide whether particular systems, sectors, or risks warrant new rules. A court need not know whether a machine has a mind before it can identify the legal subjects, select the relevant function, trace deployment control, and place responsibility where law can make it effective.
